\documentclass[11pt,letterpaper]{article}
\usepackage[T1]{fontenc}
\usepackage[utf8]{inputenc}
\usepackage[margin=0.85in,headheight=15pt]{geometry}
\usepackage{mathpazo,microtype,booktabs,longtable,array,enumitem,xcolor}
\usepackage{fancyhdr,lastpage}
\usepackage[colorlinks=true,linkcolor=teal!65!black,urlcolor=blue!55!black,bookmarksopen=true]{hyperref}
\hypersetup{pdftitle={A spreadsheet tool written in Lean},pdfauthor={LeanCalc project}}
\definecolor{ink}{HTML}{183546}
\definecolor{pale}{HTML}{EFF5F7}
\setlength{\parindent}{0pt}
\setlength{\parskip}{6pt}
\setlength{\emergencystretch}{3em}
\setlist{leftmargin=1.5em,itemsep=3pt,topsep=3pt}
\newcolumntype{P}[1]{>{\raggedright\arraybackslash}p{#1}}
\newcommand{\code}[1]{\texttt{\detokenize{#1}}}
\newcommand{\note}[1]{\par\smallskip\noindent\colorbox{pale}{\parbox{\dimexpr\linewidth-2\fboxsep}{#1}}\par\smallskip}
\pagestyle{fancy}\fancyhf{}
\fancyhead[L]{\small\textcolor{ink}{LeanCalc | Spreadsheet guide}}
\fancyhead[R]{\small User guide and development road map}
\fancyfoot[L]{\small 18 November 2026}
\fancyfoot[R]{\small\thepage\ / \pageref*{LastPage}}
\setcounter{tocdepth}{1}
\begin{document}
\thispagestyle{empty}
{\small\sffamily\color{teal!65!black} PRODUCTIVITY TOOLS / LEANCALC\par}
\vspace{18pt}
{\LARGE\bfseries\color{ink} A spreadsheet tool written in Lean\par}
\vspace{9pt}
{\large Sparse sheets, workbook objects, and OpenDocument interchange\par}
\vspace{8pt}
LeanCalc 0.5.1 \quad | \quad Lean 4.35.0-rc3 \quad | \quad 18 November 2026
\vspace{14pt}\hrule\vspace{12pt}
\note{\textbf{Version 0.5.1.} Multiple named sheets, cross-sheet and
anchored references, formula range picking, worksheet tabs, CSV, 81 function names,
incremental recalculation, formatting, 1,000-row sheets, and a limited ODS adapter. The engine and application commands are written in Lean, compiled to C,
and linked into a native macOS app. All \textbf{496 native/headless checks pass}.}
\vspace{10pt}
\tableofcontents
\vfill
The table of contents is clickable. Section 3 remains deliberately blank.
The user confirmed that Phase 1 navigation became quick after the grid-rendering
fix; that fix is retained. Phase 2's desktop controls still need acceptance testing.

\small Run \code{python3 lean_calc/build_document.py} to check the test catalogues
and compile this manual and the current-turn conversation record. Rebuild the app
first whenever Lean or bridge source changes.
\clearpage

\section{User guide}
\subsection{Start and try a workbook}
Open \code{dist/LeanCalc.app}. Save work and quit an older running copy first;
the new window title says \code{LeanCalc 0.5.1}. The app targets Apple silicon on
macOS 14 or later and includes its runtime. Python and Lean are build tools only.
The diagnostic launcher, \code{dist/Diagnose LeanCalc.command}, starts this exact
bundle and records startup information.

The selected column-chart artwork is now the app icon. Drag one \code{.lcalc},
\code{.csv}, or supported \code{.ods} file onto LeanCalc in Finder or the Dock.
Files dropped during startup are queued until the interface is ready. An already
running app opens the dropped file using its existing save/discard/cancel prompt.
CSV and ODS open as new unsaved workbooks; save them as \code{.lcalc} to retain
native details. Dropped CSV uses commas; semicolon import remains in the File menu.
File $>$ Open accepts the same three formats. Multiple-file drops are rejected
with a message so no workbook is silently replaced. Unsupported ODS features keep
the existing explicit diagnostics.

The implementation uses \href{https://www.tcl-lang.org/man/tcl9.1/TkCmd/tk_mac.html}{Tk's macOS OpenDocument event}
and \href{https://developer.apple.com/library/archive/documentation/General/Reference/InfoPlistKeyReference/Articles/CoreFoundationKeys.html}{Apple's document-type registration}.
The handler and startup queue are regression-tested headlessly. End-to-end Finder
and Dock gestures still need a desktop check in this environment.


\begin{enumerate}
\item Choose File $>$ Open and select
\code{lean_calc/examples/Phase2-workbook.lcalc}. It contains an Inputs sheet and
a Summary sheet demonstrating cross-sheet formulas, aggregates, IF, and an
absolute tax-rate reference.
\item Choose Inputs by clicking its tab below the grid. Change B2 to
\code{20}, then press Enter. Choose Summary to see the new totals.
\item Choose Sheet $>$ Add sheet, or click \textbf{+}. Names are unique
without regard to case, 1--31 characters, with no leading/trailing whitespace,
control characters, or \code{[]:*?/} or a backslash. Spaces and apostrophes are allowed.
\item Enter \code{=Inputs!D2} in a cell. For a sheet name with spaces, use
\code{='Annual total'!A1}. Double an apostrophe inside a quoted sheet name:
\code{='Bob''s data'!A1}.
\item Save as a native \code{.lcalc} workbook. All sheets are saved together.
The active sheet at the last explicit save is retained on reopening.
\end{enumerate}

A workbook contains 1--32 sheets, initially 100 rows and columns A through Z.
Sheet $>$ Expand to 1,000 rows grows the active sheet. Native version 3 saves
nonempty cells compactly; version 4 also stores workbook objects. Working sheets and display caches now store only populated entries. Thin gray grid lines and the
green active-cell outline remain. The long image strips introduced in 0.1.4
avoid expensive one-pixel tiling on repaint.

\textbf{Worksheet tabs.} Named tabs sit along the bottom of the grid. Click a
tab to switch sheets; the active tab has a green underline. Double-click a tab
to rename its sheet, or use the Sheet menu. The small \textbf{+} adds a sheet.
The arrow buttons scroll the tab strip when many sheets are open; switching
sheets automatically reveals the active tab. The name box, compact formula bar,
and quiet status line keep the main window uncluttered.

\clearpage
\subsection{Functions and presentation}
\textbf{Math.} There are 81 supported function names, including aliases. Names
are case-insensitive. Try \code{=SIN(PI()/2)}, \code{=ABS(-3)},
\code{=SQRT(81)}, \code{=ROUND(2.675,2)}, or \code{=AVERAGE(A1:A10)}.
\code{SAVERAGE} is an alias for \code{AVERAGE}. Dragging ranges still works.

\begingroup\small
\begin{longtable}{@{}P{0.20\linewidth}P{0.74\linewidth}@{}}
\toprule Family & Supported names\\\midrule
Arithmetic & ABS, SIGN, INT, TRUNC, MOD, QUOTIENT, POWER, SQRT, SQRTPI\\
Rounding & ROUND, ROUNDUP, ROUNDDOWN, FLOOR, CEILING, EVEN, ODD\\
Trigonometry & SIN, COS, TAN, ASIN, ACOS, ATAN, ATAN2; SINH, COSH, TANH, ASINH, ACOSH, ATANH\\
Constants and logs & PI, EXP, LN, LOG, LOG10, LOG2, RADIANS, DEGREES\\
Combinatorics & FACT, COMBIN, GCD, LCM\\
Aggregates & SUM, MIN, MAX, AVERAGE, SAVERAGE, COUNT, PRODUCT, SUMSQ, MEDIAN\\
Statistics & VAR, VAR.S, VAR.P, STDEV, STDEV.S, STDEV.P\\
Logic & IF, IFERROR, AND, OR, NOT\\
Criteria & COUNTIF, SUMIF, AVERAGEIF, SUMPRODUCT\\
Lookups & INDEX, MATCH, VLOOKUP (exact matches)\\
Text & LEN, LEFT, RIGHT, MID, UPPER, LOWER, TRIM, CONCATENATE, SUBSTITUTE\\
Dates & DATE, YEAR, MONTH, DAY\\\bottomrule
\end{longtable}
\endgroup
Angles are radians. Trigonometry, logarithms, roots, and noninteger powers use
binary64 math rounded to 15 significant decimal digits. Ordinary decimal and
integer arithmetic remains exact rational. ROUND rounds halfway away from zero;
INT floors, while TRUNC moves toward zero. VAR/STDEV and their .S aliases use
sample statistics; .P uses population statistics. ATAN2 takes $(x,y)$, as in Excel.
Invalid domains and overflow return \texttt{\#NUM!}; scalar text or wrong arity
returns \texttt{\#VALUE!}. Existing IF arity errors retain \texttt{\#PARSE!}.

The operators \code{^} and postfix \texttt{\%} are supported, along with scientific
literals such as \code{1e-3}. Power associates left, and unary minus binds before
power, matching Excel. Aggregate ranges skip blank/text cells. Scalar functions
such as ABS reject a rectangle; use a single cell or aggregate it first.
This is a declared function set, not a claim of complete Excel formula coverage.
Open \code{lean_calc/examples/Math-gallery.lcalc} for 61 example formulas.

\textbf{Data menu.} Find (Cmd+F) and Find next (Cmd+G) search the active sheet's
raw and displayed text, case-insensitively, wrapping at the end. Select at least
two data rows, excluding headers, to sort. The active cell's column supplies the
sort key. Equal keys keep their order, and blank cells sort last. Sorting moves
only the selected rectangle and is one undo step; moved formulas adjust relative
row references. References outside the rectangle continue to refer to coordinates.
Filter matches displayed text in the active column. Clear filter restores all rows.
Filtering is a temporary view and does not remove data from exports.

\textbf{Format menu.} Choose General, Number, Currency, Percent, Fraction, Date, Bold,
a font family/size, column alignment, or column width. Number/currency/percent
use two decimal places and never change the stored value. Currency currently
uses a dollar symbol. Fonts, formats, and explicit column settings persist and
are undoable. Mouse-dragged widths remain temporary; use Format $>$ Column width
to save a width. Row height is still fixed, so larger fonts may need desktop review.

\textbf{Print export.} File $>$ Export printable sheet saves an HTML file of the
active sheet's displayed values. Open it in a browser and choose Print or Save
as PDF. It includes a landscape print layout and escaped workbook text.

\textbf{Measured scale.} A 1,000-row by 26-column fixture with 500 formulas
referencing A1 measured 4 ms initial display, 4 ms after editing A1, and 40 MiB
peak resident memory on this host. Only 501 cells were recalculated. This native
headless benchmark excludes desktop events and repainting. The target for this
fixture is under 250 ms per edit and under 256 MiB peak memory. Forty edit
sequences are checked against a full-calculation oracle.

\subsection{OpenDocument subset: Phase 4 preview}
File $>$ Import ODS reads supported values/formulas/styles into a new unsaved
workbook. File $>$ Export ODS writes all sheets. The original imported file is
not the native Save target. Save \code{.lcalc} to retain exact rational values
and all native details. The supplied \code{lean_calc/examples/Interoperability.ods}
exercises two sheets, references, math, Unicode, and basic styles.

Supported interchange includes numeric/text/boolean cells, the implemented
formula set and references, two-decimal number/currency/percent formats,
fraction styles, bold, font family/size, and column widths/alignment.
Numeric values export at 15 significant digits; raw numeric spelling and exact
rational identity are intentional conversions. Fonts may substitute to
Helvetica, Times, or Menlo. ODS does not preserve LeanCalc undo history.

Merged cells, unsupported formulas, detected drawings, validation, named
expressions, annotations, and unsupported style properties are rejected.
Complex layout and styles defined only in a separate \code{styles.xml} are
outside this initial subset. The reader limits XML size and repeated cells and
rejects DTD/entity declarations. The app does not support arbitrary ODS files.
The package and formula conversion follow the
\href{https://docs.oasis-open.org/office/OpenDocument/v1.3/OpenDocument-v1.3-part2-packages.html}{OASIS package specification} and
\href{https://docs.oasis-open.org/office/OpenDocument/v1.3/OpenDocument-v1.3-part3-schema.html}{ODF schema}.

\textbf{Acceptance status.} Generated ODS round trips and validation against the OASIS ODF 1.3 XML schema
pass, but comparison with
Apache OpenOffice/LibreOffice has not been performed here. Phase 4 is a preview;
XLSX and executable macros are not implemented. Native Phase 5 features are
described below; their desktop acceptance remains pending.
Desktop appearance/dialog acceptance also remains pending because Tk cannot
start inside this sandbox. These gates are not counted as passing tests.

\subsection{Workbook objects and advanced functions}
Select a rectangle before using the Insert or Data menus. Save as \code{.lcalc}
to retain objects. Open \code{lean_calc/examples/Features-gallery.lcalc} for names,
a custom function, a conditional format, a chart and an embedded PNG; open
\code{Pivot-gallery.lcalc} for a refreshable summary. Insert $>$ Sheet objects
shows the current sheet's objects in a separate tabbed window. Insert $>$ Remove
object removes the numbered definition; Undo restores it. Chart views refresh
when reopened. These windows keep the main cell grid uncluttered.

\begin{longtable}{@{}P{0.22\linewidth}P{0.72\linewidth}@{}}
\toprule Feature & Operation and declared limits\\\midrule
Named ranges & Insert $>$ Named range. Names use letters/underscores, are globally
unique without case, and cannot shadow built-ins. Use \code{=SUM(Revenue)} or a
single-cell name in arithmetic. Names track sheet renames, row/column edits and
complete-range cuts. A cut splitting a named range is rejected.\\
Validation & Data $>$ Validation. Enter \code{number,0,100}, \code{whole,0,100}, or
\code{list,Apple,pear}. Blanks are allowed; lists match exact displayed strings.
Invalid edits, dependent recalculations, paste, fill, CSV replacement and sort are
rejected atomically. Escape cancels an invalid pending draft. Removing a rule is undoable.\\
Conditional formatting & Data $>$ Conditional format. Enter \code{>,15,green}.
Comparisons are \code{< > = <= >= <>}; colors are red, green, yellow or blue.
Rules use evaluated numbers, skip blanks/text, and the last matching rule wins.
Colors appear in the grid and printable HTML.\\
Charts & Insert $>$ Chart; choose bar or line. Up to 200 selected rows; first
column supplies labels and second supplies values, or one column supplies values
with row-number labels. Blank/text values are skipped, errors reported. Charts
appear in the objects window and as SVG in printable HTML.\\
Images & Insert $>$ PNG image embeds the file, anchored at the selected cell.
Maximum 1.4 MB compressed, 4 million pixels, 4,096 pixels per side. The objects
window scales the preview; printable HTML embeds it. Native save retains the
original PNG; no external file is needed on reopening. A workbook whose
embedded PNG is truncated or corrupted is refused on opening.\\
Pivot tables & Select a range including one header row; Insert $>$ Pivot table.
Choose a new sheet name and, for example, \code{A,B,SUM}. One grouping column and
one value column; SUM, COUNT (nonempty values), AVERAGE, MIN or MAX. Blank groups
are skipped; numeric aggregates ignore text. Data $>$ Refresh pivot table updates
the active summary in one undo step and replaces its output cells. Source sheet
renames and structural edits update its definition.\\
Extensibility & Insert $>$ Custom function: \code{DOUBLE(X)=X*2}. Parameters are
letters/underscores, at most 16; functions may use other formulas/names. The usual
length/depth/work limits apply, including to recursive definitions. Data exports
and imports versioned JSON function packs. They contain expressions, with no
shell, file, network, plugin-loading or macro execution.\\\bottomrule
\end{longtable}
Up to 128 objects per sheet. Copy/fill copy cell values and basic cell styles;
they do not duplicate globally named objects. Range objects stay at coordinates
when sorting. Complete cuts move local objects; partial cuts that would split
objects are rejected. Deleting an entire named range removes its definition.
Pivot headers and key/value columns cannot be structurally deleted while the
pivot definition exists. Deleting its source sheet leaves a definition whose
refresh reports the missing source. Preview windows show a snapshot; reopen to
see edits. The native workbook remains the source of truth.

\textbf{Conditional aggregates and lookup.} COUNTIF, SUMIF and AVERAGEIF accept
one rectangular criterion range and an optional matching sum/average range.
Numeric comparisons and case-insensitive literal text are supported; wildcards
are literal characters. SUMPRODUCT accepts two equally shaped ranges and ignores
text. INDEX uses one-based row/column positions. MATCH requires an explicit zero
match mode and a single row or column. VLOOKUP requires the fourth argument
FALSE/zero; approximate modes return \texttt{\#N/A}. Missing exact matches return
\texttt{\#N/A}; invalid positions return \texttt{\#REF!}.

\textbf{Text and dates.} Text functions count Unicode code points; LEFT/RIGHT
without a count return one character. MID uses a one-based start. SUBSTITUTE
replaces all literal occurrences; its optional occurrence selector is outside
this subset. TRIM collapses ordinary spaces. Dates use the Gregorian calendar,
an epoch of 1899-12-30, and the supported interval 1900-01-01 through 9999-12-31.
DATE requires valid year/month/day integers and rejects rollover dates and the
fictitious 1900-02-29. YEAR/MONTH/DAY require integral serials. Format $>$ Date
shows YYYY-MM-DD without changing the stored number. These explicit limits are
narrower than Excel's coercion and rollover behavior. The epoch follows the
\href{https://docs.oasis-open.org/office/OpenDocument/v1.3/OpenDocument-v1.3-part4-formula.html}{OASIS OpenFormula date model}.

\textbf{Interchange policy.} ODS export refuses workbooks with native objects or
date formatting until those objects/styles have a tested ODS mapping. It does not
silently discard them. CSV remains an explicit values-only export. Page setup,
view preferences and application metadata are outside the declared ODS subset;
unsupported workbook content is rejected and the original file is never replaced.
See \code{lean_calc/phase-completion.txt} for the remaining acceptance gates.

\subsection{Editing and rectangular selection}
Click a cell and type to replace it. Double-click, click the active cell, or press
F2 to edit its existing value with a caret. The formula bar and cell editor share
one pending draft. Enter commits and moves down; Shift+Enter moves up. Tab and
Shift+Tab commit and move forward/backward, including across row boundaries.
Apply commits without moving. Escape cancels the draft. Arrows navigate cells
when not editing and move the caret while editing.

Use Shift-click, Shift-arrow keys, or drag from a newly clicked cell to extend a
rectangular selection. The selected rectangle is shaded and its address appears
at the top. A plain click or arrow returns to one cell. Typing starts a new edit
in the active cell and clears the rectangular selection. The cell editor follows
scrolling and hides when its cell is outside the viewport.

\textbf{Pick a range while entering a formula.} In an empty result cell, type
\code{=sum(}. Drag across a rectangle of numbers on the current sheet, release
the mouse, type \code{)}, and press Return. For example, dragging A1 through B3
inserts \code{A1:B3}, so the result cell receives \code{=sum(A1:B3)}. The referenced
rectangle is highlighted while the draft stays in the original cell. Dragging
changes only the draft; Return or Apply commits it, and Escape cancels it.

This also works from the formula bar, after an operator or comma, and before an
existing closing parenthesis. Selecting an existing reference inside the formula
lets a drag replace that text. A single picked cell inserts just its address.
Dragging in the opposite direction gives the same normalized rectangle. After
release, continue typing to finish the expression or add another argument.

\textbf{Copy and paste.} Cmd+C copies the selected rectangle. Select the top-left
destination and press Cmd+V. A selection at the destination uses its top-left
corner. A copy within this app carries the original coordinates, so relative
references adjust. External tab-separated clipboard text is parsed as a table;
formulas in it are used as supplied. Quoted fields preserve tabs, newlines, and
quotes. Uneven rows and out-of-bounds pastes are rejected before any pasted data
is written. A successful paste is one undo step.

\textbf{Cut and paste.} Cmd+X prepares a move; source cells remain until Paste.
The complete move is one undo step, including any references updated elsewhere.
Changing workbook content after Cut invalidates that pending move; select and Cut
again. Merely switching the active sheet does not invalidate it. Overlapping
moves use the original cells. A move that would split a referenced range is
rejected: copy instead, or move the entire referenced range.

\textbf{Fill.} Select a rectangle and use Edit $>$ Fill Down (Cmd+D) to copy its
first row downward, or Fill Right (Cmd+R) to copy its first column rightward.
Formulas adjust just as on Copy. Edit $>$ Clear cells, Delete, or Backspace clears
the rectangle. Each operation is undoable as a whole.

\subsection{References, structural edits, and undo}
Reference behavior was specified before implementation in
\code{lean_calc/phase2-reference-rules.txt}.

\begin{center}\small
\begin{tabular}{@{}lll@{}}
\toprule\textbf{Reference} & \textbf{Copy behavior} & \textbf{Example: one row/right}\\\midrule
\code{A1} & Row and column move & \code{B2}\\
\code{$A$1} & Both stay fixed & \code{$A$1}\\
\code{A$1} & Column moves & \code{B$1}\\
\code{$A1} & Row moves & \code{$A2}\\
\code{Other!A1} & Same sheet; both coordinates move & \code{Other!B2}\\
\bottomrule
\end{tabular}\end{center}
Unqualified references use the formula's destination sheet after Copy. Explicit
sheet qualifiers stay explicit. References shifted outside A1:Z100 become
\texttt{\#REF!}. Rewriting may normalize capitalization, remove formula whitespace,
and quote sheet names; quoted string contents are preserved.

Cut moves cell identities: references to moved cells follow them, even if they
have dollar anchors. A moved formula retains its references to cells outside the
block; a cross-sheet move adds the original sheet qualifier when needed. The
range-splitting restriction above prevents an ambiguous partial range move.

The Sheet menu inserts or deletes one row/column at the selection's top/left
edge. Every sheet's references follow the structural change, regardless of dollar
anchors. Deleted single-cell targets become \texttt{\#REF!}. Ranges shrink on deletion,
expand when insertion is inside the range, and become \texttt{\#REF!} if no cells
remain. Insertion before the first range cell shifts the whole range. The grid
size stays fixed: insertion is refused if it would discard nonempty edge cells;
deletion adds an empty row or column at the far edge.

Renaming a sheet rewrites qualified references throughout the workbook. Deleting
a sheet invalidates references to it; at least one sheet must remain. Both can be
undone. Cmd+Z and Shift+Cmd+Z undo/redo workbook changes across sheets and restore
the relevant active sheet. Up to 100 snapshots are retained. Navigation alone is
not an undo step. Editing after undo clears redo history.

\subsection{Values and formulas}
Cells hold text, decimal numbers, or formulas beginning with \code{=}. A leading
apostrophe makes literal text: \code{'=1+2} displays \code{=1+2}. Cells have a
4096-character limit and cannot contain NUL. Arithmetic uses exact rational
numbers; \code{0.1+0.2} is exactly \code{0.3}. Terminating results with up to ten
decimal places display as decimals; other results display as fractions.

Arithmetic supports \code{+ - * /}, unary signs and parentheses. Comparisons
\code{= <> < <= > >=} return 1 or 0; arithmetic has higher precedence. Equality
can compare text; ordered comparisons require numeric values. TRUE and FALSE
are 1 and 0. A blank reference is zero in arithmetic; text is a value error.

\begin{center}\small
\begin{tabular}{@{}P{0.25\linewidth}P{0.66\linewidth}@{}}
\toprule\textbf{Function} & \textbf{Behavior}\\\midrule
\code{SUM(...)} & Add scalar expressions and rectangular ranges.\\
\code{MIN(...)} / \code{MAX(...)} & Smallest/largest numeric value; zero if none.\\
\code{AVERAGE(...)} & Numeric mean; division error if no numeric values.\\
\code{COUNT(...)} & Number of numeric values; blank/text references are skipped.\\
\code{IF(test,yes,no)} & Exactly three arguments; numeric zero chooses no. Only the chosen branch is evaluated. Branches may return text.\\
\bottomrule
\end{tabular}\end{center}
Aggregates skip blank/text cells and propagate formula errors. They accept multiple
arguments, for example \code{=SUM(A1:B3,5,Other!C1)}. Range endpoints can be reversed;
in \code{Other!A1:B3}, the second endpoint inherits Other. A range cannot span
sheets. Example: \code{=IF(SUM(A1:A5)>100,"Over budget","OK")}. Double double-quotes
inside a string to quote a double-quote. Unsupported functions produce a name error.

\subsection{Errors and calculation limits}
\begin{center}\small
\begin{tabular}{@{}ll@{}}
\toprule\textbf{Error} & \textbf{Meaning}\\\midrule
\texttt{\#DIV/0!} & Division by zero or an average with no numeric values.\\
\texttt{\#CYCLE!} & Circular reference, including across sheets.\\
\texttt{\#VALUE!} & Text in numeric arithmetic or an unsupported range context.\\
\texttt{\#REF!} & Missing sheet, deleted target, or out-of-grid reference.\\
\texttt{\#NAME?} & Unsupported name or function.\\
\texttt{\#PARSE!} & Invalid formula syntax.\\
\texttt{\#LIMIT!} & Formula length, depth or work budget exceeded.\\
\bottomrule
\end{tabular}\end{center}
Formulas have a 512-character limit, parser/evaluation depth budgets of 128, and
20,000 evaluation steps per result. IF does not evaluate its unused branch, so an
error or cycle there does not affect the result. Recalculation after edits remains
synchronous; complex workbooks can take time to recalculate. Navigation reuses
cached display values. Changes in another sheet invalidate dependent displayed
results when that sheet is viewed. Named ranges participate in dependency tracking.
Workbooks containing custom functions conservatively clear the result cache after
an edit. Executable macros are not implemented; function packs contain expressions only.

\subsection{CSV import and export}
\textbf{Save As CSV.} Choose File $>$ Save As and enter a filename ending in
\code{.csv}. LeanCalc writes UTF-8 comma-separated values, with CSV quoting,
instead of the native workbook format. Both \code{.csv} and \code{.lcalc} are
accepted by the save dialog. The \code{.CSV} extension works too.

CSV contains only the active sheet's displayed values. Exporting leaves the
native workbook path and unsaved-change indicator intact. Save to \code{.lcalc}
to retain formulas and all sheets. If CSV is chosen while saving before closing,
the export completes and the workbook stays open. File $>$ Export CSV remains
available, including the semicolon option.

The File menu offers separate comma and semicolon import/export commands. Files
must be UTF-8; an optional UTF-8 BOM is accepted on import. Fields use double-quote
escaping, doubled double-quotes, and may contain delimiters or CR/LF inside quotes.
Import accepts CRLF, LF, or CR record endings. Export uses CRLF.

Import validates the entire file, then asks to replace the active sheet. Other
sheets and the native save path are preserved. Undo restores the previous sheet.
Shorter rows are padded with empty cells. More rows than the active sheet capacity (100 or expanded to 1,000), more than 26 columns,
malformed quotes, NUL, and oversized fields are rejected. Imported formula-looking
text is stored literally, never executed; an internal leading apostrophe preserves
that distinction. An already literal apostrophe is also preserved on export.

Export writes \emph{displayed values}, not formulas, from the smallest A1-anchored
rectangle containing nonempty raw cells. A blank sheet produces one empty field.
It does not change the native save path or mark the workbook saved. Use native
\code{.lcalc} to retain formulas and multiple sheets. A limited ODS adapter is available; XLSX is not yet supported.

\subsection{Native files and migration}
\code{LeanCalc-2} files store all sheet names, raw cell contents, and the active
sheet as UTF-8 JSON. Existing \code{LeanCalc-1} files open as a workbook containing
Sheet1. Formatting or expanding a sheet writes version 3. Workbook objects write version 4.
This build reads versions 1--4; older builds cannot read version 4.
Use Save As if a separate old-format original is needed. Undo history, pending
clipboard moves, and the cell selection are session state.

Save commits a pending draft and compares the file with the bytes last read or
saved. Conflicts are reported instead of overwritten. Save As can preserve both
versions. The C bridge writes a temporary file beside the target, synchronizes it,
and atomically renames it. This is not a file lock; a concurrent writer can race
the comparison and rename. Native Save resolves symlinks before replacement.

\clearpage
\subsection{Menus and shortcuts}
\begingroup\small
\begin{longtable}{@{}P{0.30\linewidth}P{0.64\linewidth}@{}}
\toprule\textbf{Action} & \textbf{Behavior / shortcut}\\\midrule\endhead
File: New / Open & New workbook / open .lcalc. Cmd+N / Cmd+O; offer to save changes.\\
File: Save / Save As & Save workbook as .lcalc; .csv exports active sheet values. Cmd+S / Shift+Cmd+S.\\
File: Import CSV & Replace active sheet after confirmation; choose comma or semicolon.\\
File: Export CSV & Export active sheet's displayed values; choose comma or semicolon.\\
File: Quit & Save/discard/cancel prompt. Cmd+Q.\\
Edit: Undo / Redo & Workbook-wide history. Cmd+Z / Shift+Cmd+Z.\\
Edit: Cut / Copy / Paste cells & Rectangular clipboard. Cmd+X / Cmd+C / Cmd+V.\\
Edit: Clear cells & Clear the selection. Delete or Backspace when navigating.\\
Edit: Fill Down / Right & Repeat first row/column with adjusted references. Cmd+D / Cmd+R.\\
Sheet: Add / Rename / Delete & Manage named sheets; all undoable. The + button adds a sheet.\\
Sheet: Insert/Delete row/column & One structural change at the selection's top/left edge.\\
Insert & Names, charts, PNG images, pivot tables, custom functions and object management.\\
Data & Validation, conditional formats, pivot refresh, function packs, find, sort and filter.\\
Format & Numeric/date display, fonts, column widths and alignment.\\
File: Print / ODS & Printable HTML; import/export the declared ODS subset.\\
Help & Compact guide to operations and references.\\
\bottomrule
\end{longtable}\endgroup
Command shortcuts also have Control aliases. Clipboard shortcuts act on selected
text while editing in the formula bar or cell editor, and on cells while navigating.
The Edit menu's cell clipboard commands always operate on cells. Shift-click and
Shift-arrows extend the selection; Enter/Tab commit and move; Escape cancels; F2 edits.

\subsection{Build and validation}
The compatibility engine is in \code{LeanCalcCore.lean}; the new workbook engine
is in \code{Phase2.lean}. Tests are in \code{Phase2Tests.lean}, and the desktop
application and integration checks are in \code{LeanCalc.lean}. The builder compiles
the Lean modules to C and links them with the Tcl/Tk, filesystem, ZIP and XML
bridges. The installed macOS XML/archive libraries and zip utility handle ODS.
It regenerates both PDFs, signs the app locally, and packages the ZIP.

\begin{quote}\small\ttfamily
python3 lean\_calc/build.py\\
dist/LeanCalc.app/Contents/MacOS/LeanCalc --gui-self-test\\
python3 lean\_calc/build\_document.py
\end{quote}
The toolchain is Lean 4.35.0-rc3, Apple command-line tools, Python 3.9+, Homebrew
static GMP/libuv, Tcl/Tk 9.0.4, and LaTeX/latexmk. The supplied runtime path can be
overridden. Keep generated C files unedited.

\textbf{Status.} All 496 native/headless checks pass and the app signature verifies.
The user accepted Phase 1's responsiveness. The Phase 2 desktop suite still aborts
at macOS GUI initialization in this sandbox, so new controls, dialogs, visual
selection, and Finder activation need desktop acceptance. Headless tests exercise
real Lean handlers and Tcl translation, but substitute widget/dialog adapters.
The navigation benchmark still performs zero grid-row writes for unchanged cells;
it excludes painting. \code{--gui-render-benchmark} measures real Tk keyboard and
painting work on a desktop and writes \code{dist/LeanCalc-render-benchmark.log}.

\clearpage
\section{Unit tests}
There are \textbf{496 named native/headless checks}: the original 61 baseline
checks, 17 editing-interaction checks, 12 display-cache checks, 87 Phase 2 engine
checks, 23 Phase 2 command-handler checks, 25 formula-picking checks, 19 saving/tab
checks, 82 math checks, 30 Phase 3 checks, 15 ODS checks, 8 Phase 3 interaction checks, 87 Phase 5 checks, and 13 Phase 5
interaction checks, and 17 document-open checks. A named check can contain several
assertions. One differential check compares the new engine with the original
arithmetic engine across 308 inputs. These are finite tests, not formal proofs.

The document builder verifies the baseline names against \code{tests.json} and
regenerates the catalogues below from executable results. The separate desktop
GUI suite is excluded from the passing count. Original checks retain their names;
legacy single-sheet encoding tests cover the Phase 1 migration input format.
Workbook-format and migration tests cover the new active format.

\subsection{Original 61-check baseline}
\begin{longtable}{@{}P{0.07\linewidth}P{0.87\linewidth}@{}}
\toprule\textbf{No.} & \textbf{What the test checks}\\\midrule\endfirsthead
\toprule\textbf{No.} & \textbf{What the test checks (continued)}\\\midrule\endhead
% BEGIN TEST CATALOGUE
01 & First cell address.\\
02 & Last cell address.\\
03 & Lowercase addresses.\\
04 & Whitespace around addresses.\\
05 & Reject row zero.\\
06 & Reject row beyond sheet.\\
07 & Reject column beyond sheet.\\
08 & Reject missing row.\\
09 & Address round-trip across all 2600 cells.\\
10 & Empty sheet dimensions and contents.\\
11 & Parse positive integer.\\
12 & Parse negative decimal.\\
13 & Parse leading decimal point.\\
14 & Parse explicit plus.\\
15 & Reject multiple decimal points.\\
16 & Reject lone decimal point.\\
17 & Exact decimal addition.\\
18 & Display terminating decimal.\\
19 & Display repeating result as fraction.\\
20 & Preserve Unicode text.\\
21 & Formula addition.\\
22 & Formula subtraction.\\
23 & Formula multiplication.\\
24 & Formula division.\\
25 & Multiplication precedence.\\
26 & Parentheses precedence.\\
27 & Left-associative subtraction.\\
28 & Left-associative division.\\
29 & Unary minus and plus.\\
30 & Formula whitespace.\\
31 & Reference another cell.\\
32 & Recalculate changed dependency.\\
33 & Blank references equal zero.\\
34 & Text reference is an error.\\
35 & Division by zero error.\\
36 & Direct circular reference.\\
37 & Indirect circular reference.\\
38 & Out-of-range reference error.\\
39 & Unknown function error.\\
40 & Reject trailing expression tokens.\\
41 & Reject unmatched parentheses.\\
42 & Rectangular SUM.\\
43 & Reversed SUM endpoints.\\
44 & SUM skips text and blank cells.\\
45 & SUM propagates formula errors.\\
46 & Formula length limit.\\
47 & Apostrophe forces literal text.\\
48 & Editing affects only chosen cell.\\
49 & Out-of-bounds edit is unchanged.\\
50 & Undo restores previous sheet.\\
51 & Redo restores undone edit.\\
52 & New edit clears redo history.\\
53 & Identical edit does not create undo history.\\
54 & File round-trip preserves formulas and Unicode.\\
55 & Reject malformed file.\\
56 & Reject unsupported file version.\\
57 & Reject incorrect sheet dimensions.\\
58 & Reject oversized or NUL-containing cells.\\
59 & Native main thread and noninteractive startup.\\
60 & Tcl bridge preserves text without evaluation.\\
61 & Atomic save round-trip and conflict protection.\\
% END TEST CATALOGUE
\bottomrule\end{longtable}

\subsection{Existing interaction and cache checks}
The 17 interaction checks exercise click/keyboard selection, draft editing,
commit/cancel, undo/redo, formula source editing, bounds, Unicode, shortcut
filtering, validation and saving a pending draft. The 12 cache checks cover initial
painting requests, unchanged navigation, changed rows, dependent recalculation,
formula selection, undo/redo, equivalent formulas, clearing, errors, file loading,
and New. They continue to pass with the workbook model.

\subsection{Phase 2 engine and integration checks}
Entries 1--87 cover the engine; entries 88--110 exercise real application commands
with headless widget, clipboard, and dialog adapters.
\begin{longtable}{@{}P{0.07\linewidth}P{0.87\linewidth}@{}}
\toprule\textbf{No.} & \textbf{What the test checks}\\\midrule\endfirsthead
\toprule\textbf{No.} & \textbf{What the test checks (continued)}\\\midrule\endhead
% BEGIN PHASE2 CATALOGUE
01 & Phase 2 arithmetic precedence and exact decimals.\\
02 & Unary signs and subtraction remain compatible.\\
03 & All comparison operators.\\
04 & Comparison precedence follows arithmetic.\\
05 & IF evaluates only the chosen branch.\\
06 & IF returns text and accepts nested conditions.\\
07 & Quoted strings preserve punctuation and escaped quotes.\\
08 & IF rejects missing or extra arguments.\\
09 & MIN MAX AVERAGE COUNT and variadic SUM.\\
10 & Empty aggregates have defined results.\\
11 & Aggregate errors propagate.\\
12 & Unknown functions and malformed argument lists fail.\\
13 & Range functions skip blank and literal text cells.\\
14 & Cross-sheet formulas calculate.\\
15 & Quoted sheet names and absolute references calculate.\\
16 & Range endpoints inherit the first sheet qualifier.\\
17 & Sheet and cell lookup is case-insensitive.\\
18 & Unknown sheets report REF.\\
19 & Cross-sheet cycles are detected.\\
20 & Unchosen IF branches do not cause cross-sheet cycles.\\
21 & Formula length budget remains enforced.\\
22 & Relative absolute and mixed copy adjustments.\\
23 & Copy adjusts qualified references.\\
24 & Copy preserves formula string literals.\\
25 & Out-of-bounds relative copies become REF.\\
26 & Range copies preserve endpoint anchors.\\
27 & Adding a sheet activates it.\\
28 & Duplicate names are rejected case-insensitively.\\
29 & Invalid sheet names are rejected.\\
30 & Quoted sheet names escape apostrophes.\\
31 & Renaming updates references in every sheet.\\
32 & Renaming to another sheet name is rejected.\\
33 & Deleting a sheet invalidates its references.\\
34 & Last sheet cannot be deleted.\\
35 & Undo crosses sheets and restores active sheet.\\
36 & Redo restores workbook operations.\\
37 & Editing after undo clears workbook redo.\\
38 & Rectangular paste fills all cells.\\
39 & Rectangle copy adjusts relative formulas.\\
40 & Paste rejects overflow and ragged tables.\\
41 & Paste is one undo step.\\
42 & Fill Down copies the first row and adjusts mixed references.\\
43 & Fill Right adjusts columns and preserves absolute rows.\\
44 & Clear rectangle preserves outside cells.\\
45 & Overlapping Cut uses the original rectangle.\\
46 & Cut across sheets updates outside and anchored references.\\
47 & Moved formula retains references outside its block.\\
48 & Cut move is one undo operation.\\
49 & Insert row shifts data without truncation.\\
50 & Insert row expands a spanning range.\\
51 & Structural changes shift absolute cross-sheet references.\\
52 & Delete row shrinks ranges and shifts data.\\
53 & Deleted single-cell reference becomes REF.\\
54 & Deleting a range boundary retains surviving cells.\\
55 & Reversed ranges shrink correctly.\\
56 & Deleting a whole referenced range becomes REF.\\
57 & Insert column shifts formulas and their references.\\
58 & Delete column reverses a prior empty insertion.\\
59 & Insertion cannot discard nonempty edge cells.\\
60 & Structural operation is one undo step.\\
61 & Phase 1 files migrate without losing cells.\\
62 & Phase 2 round-trip preserves sheets and active sheet.\\
63 & Version and active-sheet validation.\\
64 & Workbook format rejects duplicate names and invalid shapes.\\
65 & CSV handles BOM quoted delimiters doubled quotes and newlines.\\
66 & CSV round-trip preserves embedded CR LF quotes and Unicode.\\
67 & Semicolon delimiter is explicit.\\
68 & CSV accepts CR LF and CRLF record endings.\\
69 & CSV preserves trailing empty fields.\\
70 & CSV rejects malformed quoting.\\
71 & CSV rejects too many rows or columns.\\
72 & CSV rejects NUL and oversized fields.\\
73 & CSV import pads shorter rows and keeps formulas literal.\\
74 & CSV import preserves other sheets.\\
75 & CSV import is one undo step.\\
76 & CSV export writes displayed values and CRLF.\\
77 & Blank sheet CSV exports one empty cell.\\
78 & CSV literal apostrophes survive import and export.\\
79 & Clipboard TSV round-trip preserves tabs newlines and quotes.\\
80 & COUNT ignores a blank scalar reference.\\
81 & Quoted sheet names can contain exclamation marks.\\
82 & Cut refuses to split a referenced range.\\
83 & Deleting the first range row preserves a surviving formula.\\
84 & Deleting the last range row preserves a surviving formula.\\
85 & Workbook enforces its 32-sheet limit.\\
86 & Workbook undo history remains bounded.\\
87 & Phase 2 agrees with the Phase 1 arithmetic oracle across 308 cases.\\
88 & Shift navigation selects a rectangle.\\
89 & Copy does not alter workbook or undo history.\\
90 & Clipboard paste shifts relative formulas in the real handler.\\
91 & UI undo reverses rectangular paste as one action.\\
92 & UI redo restores the pasted rectangle.\\
93 & Fill Down handler updates the whole selected rectangle.\\
94 & New sheet and cross-sheet formula are connected to the UI.\\
95 & Sheet rename updates live cross-sheet formulas.\\
96 & Cut leaves source cells intact until Paste.\\
97 & Cut across sheets moves data and follows incoming references.\\
98 & UI undo restores a cut across sheets.\\
99 & Stale Cut is rejected without losing data.\\
100 & External TSV clipboard pastes as a rectangle.\\
101 & Row insertion handler shifts pasted data.\\
102 & Row deletion handler removes the inserted row.\\
103 & Delete sheet command remains undoable.\\
104 & Undo restores a deleted sheet and activates it.\\
105 & Native save and load preserve every sheet and the active tab.\\
106 & Workbook save still detects external file changes.\\
107 & CSV import handler replaces only the active sheet.\\
108 & CSV export writes display text without changing the native path.\\
109 & Undo reverses CSV import in one step.\\
110 & Cached results refresh after editing a referenced sheet.\\
% END PHASE2 CATALOGUE
\bottomrule\end{longtable}

\subsection{Formula range-picking checks}
These checks exercise the pointer handlers and draft editing through headless
widget adapters, including no premature commit and no grid-value rewrites.
The desktop suite separately exercises real pointer and key events.
\begin{longtable}{@{}P{0.07\linewidth}P{0.87\linewidth}@{}}
\toprule\textbf{No.} & \textbf{What the test checks}\\\midrule\endfirsthead
\toprule\textbf{No.} & \textbf{What the test checks (continued)}\\\midrule\endhead
% BEGIN PICKING CATALOGUE
01 & Picking starts without committing or moving the formula cell.\\
02 & Dragging replaces the temporary reference with a rectangle.\\
03 & Repeated drag motion never duplicates reference text.\\
04 & Mouse release leaves the draft in its original editor.\\
05 & Formula picking does not rewrite grid values.\\
06 & Closing parenthesis and Return calculate the picked SUM.\\
07 & Committing clears reference highlights and picking state.\\
08 & Undo reverses the complete picked formula in one step.\\
09 & Redo restores the picked formula.\\
10 & Reverse-direction dragging produces a normalized range.\\
11 & Picking again replaces the previous range until typing resumes.\\
12 & Escape cancels picking without changing committed content.\\
13 & Picking at the caret preserves the existing formula suffix.\\
14 & Ordinary clicking after a completed formula still commits and navigates.\\
15 & Picking replaces a selected reference inside an existing formula.\\
16 & Typing a comma permits picking another function argument.\\
17 & Formula-bar picking keeps focus and synchronizes both editors.\\
18 & Formula-bar picked ranges use the normal commit path.\\
19 & A single picked cell inserts one address without a colon.\\
20 & Picking after an arithmetic operator appends a new reference.\\
21 & Dragging outside the grid retains the last valid reference.\\
22 & Quoted strings and quoted sheet names are not insertion slots.\\
23 & Unicode and escaped quotes preserve caret insertion boundaries.\\
24 & Non-formula editing still commits normally on a cell click.\\
25 & Oversized range insertion preserves the draft and original cell.\\
% END PICKING CATALOGUE
\bottomrule\end{longtable}

\subsection{Saving and worksheet-tab checks}
The saving checks inspect files written by the real handlers. Tab checks run
the tab-control script with headless widget adapters; geometry and appearance
also have desktop checks, excluded from the passing count.
\begin{longtable}{@{}P{0.07\linewidth}P{0.87\linewidth}@{}}
\toprule\textbf{No.} & \textbf{What the test checks}\\\midrule\endfirsthead
\toprule\textbf{No.} & \textbf{What the test checks (continued)}\\\midrule\endhead
% BEGIN SAVE TABS CATALOGUE
01 & Save As with a CSV suffix writes quoted UTF-8 CSV values.\\
02 & CSV Save As is not mistaken for saving the entire workbook.\\
03 & Save dialog permits both native and CSV filename extensions.\\
04 & Native Save As still persists all sheets and marks the workbook saved.\\
05 & Uppercase CSV extension also exports and commits the pending draft.\\
06 & CSV export preserves the native path and conflict baseline.\\
07 & CSV export preserves formulas, other sheets, and unsaved changes.\\
08 & Save after CSV export still updates the native workbook.\\
09 & Cancelling Save As leaves existing files unchanged.\\
10 & CSV during save-before-close keeps the unsaved workbook open.\\
11 & CSV Save As exports the active sheet rather than the first sheet.\\
12 & A workbook renders one visible named sheet tab.\\
13 & Adding a sheet creates its named tab and marks it active.\\
14 & Clicking a tab commits a pending draft and switches sheets.\\
15 & Switching tabs removes the old active underline.\\
16 & Cell navigation leaves unchanged sheet tabs untouched.\\
17 & Renaming a sheet updates its tab with literal Unicode and punctuation.\\
18 & Deleting a sheet removes its tab and selects the remaining sheet.\\
19 & Undo restores deleted sheet tabs and the active indicator.\\
% END SAVE TABS CATALOGUE
\bottomrule\end{longtable}

\subsection{Math, Phase 3, and ODS checks}
\begin{longtable}{@{}P{0.07\linewidth}P{0.87\linewidth}@{}}
\toprule\textbf{No.} & \textbf{What the test checks}\\\midrule\endfirsthead
\toprule\textbf{No.} & \textbf{What the test checks (continued)}\\\midrule\endhead
% BEGIN ADVANCED CATALOGUE
01 & Math exact: ABS(-3.25).\\
02 & Math exact: SIGN(-2).\\
03 & Math exact: INT(-2.3).\\
04 & Math exact: TRUNC(-2.9).\\
05 & Math exact: ROUND(2.675,2).\\
06 & Math exact: ROUND(-1.5,0).\\
07 & Math exact: ROUND(1234,-2).\\
08 & Math exact: ROUNDUP(-1.21,1).\\
09 & Math exact: ROUNDDOWN(-1.29,1).\\
10 & Math exact: TRUNC(12.345,2).\\
11 & Math exact: MOD(-3,2).\\
12 & Math exact: MOD(3,-2).\\
13 & Math exact: QUOTIENT(-7,2).\\
14 & Math exact: POWER(2,10).\\
15 & Math exact: 2\textasciicircum{}-3.\\
16 & Math exact: 2\textasciicircum{}3\textasciicircum{}2.\\
17 & Math exact: -2\textasciicircum{}2.\\
18 & Math exact: 50\%*8.\\
19 & Math exact: FACT(5).\\
20 & Math exact: COMBIN(6,2).\\
21 & Math exact: GCD(12,18,24).\\
22 & Math exact: LCM(4,6).\\
23 & Math exact: FLOOR(7.8,2).\\
24 & Math exact: CEILING(7.1,2).\\
25 & Math exact: EVEN(-3).\\
26 & Math exact: ODD(0).\\
27 & Math exact: PRODUCT(2,3,4).\\
28 & Math exact: SUMSQ(2,3).\\
29 & Math exact: MEDIAN(5,1,3).\\
30 & Math exact: MEDIAN(4,1,2,3).\\
31 & Math exact: VAR.S(1,2,3).\\
32 & Math exact: VAR.P(1,2,3).\\
33 & Math exact: STDEV.S(1,2,3).\\
34 & Math exact: SAVERAGE(2,4).\\
35 & Math exact: AND(1,2,0).\\
36 & Math exact: OR(0,2).\\
37 & Math exact: NOT(0).\\
38 & Math exact: IFERROR(1/0,7).\\
39 & Math exact: IFERROR(5,1/0).\\
40 & Math exact: 1e-3+2E+2.\\
41 & Math approximate: SIN(PI()/2).\\
42 & Math approximate: COS(0).\\
43 & Math approximate: TAN(0).\\
44 & Math approximate: ASIN(0).\\
45 & Math approximate: ACOS(1).\\
46 & Math approximate: ATAN(0).\\
47 & Math approximate: ATAN2(1,0).\\
48 & Math approximate: SINH(0).\\
49 & Math approximate: COSH(0).\\
50 & Math approximate: TANH(0).\\
51 & Math approximate: ASINH(0).\\
52 & Math approximate: ACOSH(1).\\
53 & Math approximate: ATANH(0).\\
54 & Math approximate: SQRT(81).\\
55 & Math approximate: POWER(9,0.5).\\
56 & Math approximate: EXP(0).\\
57 & Math approximate: LN(1).\\
58 & Math approximate: LOG(100).\\
59 & Math approximate: LOG(8,2).\\
60 & Math approximate: LOG10(1000).\\
61 & Math approximate: LOG2(8).\\
62 & Math approximate: DEGREES(PI()).\\
63 & Math approximate: RADIANS(180)/PI().\\
64 & Math approximate: SIN(1)\textasciicircum{}2+COS(1)\textasciicircum{}2.\\
65 & Math error: SQRT(-1).\\
66 & Math error: LN(0).\\
67 & Math error: LOG(3,1).\\
68 & Math error: ASIN(2).\\
69 & Math error: POWER(-2,0.5).\\
70 & Math error: EXP(10000).\\
71 & Math error: MOD(1,0).\\
72 & Math error: POWER(0,-1).\\
73 & Math error: VAR.S(1).\\
74 & Math error: SIN().\\
75 & Math error: ABS(1,2).\\
76 & Math error: PI(1).\\
77 & Math error: ABS("text").\\
78 & Math error: FACT(-1).\\
79 & Math error: POWER(2,1025).\\
80 & New aggregates accept ranges and ignore referenced text.\\
81 & Scalar math rejects a rectangle instead of silently flattening it.\\
82 & Scientific literal signs survive copy rewriting.\\
83 & Initial incremental display agrees with full evaluation.\\
84 & Unchanged workbook reuses every displayed result.\\
85 & An edit recalculates only its transitive dependents.\\
86 & Replacing a formula removes obsolete dependencies.\\
87 & Dependency cycles terminate and display cycle errors.\\
88 & Undoing a cycle invalidates all affected cached errors.\\
89 & Cross-sheet edits invalidate cached dependents on another sheet.\\
90 & IF branch dependencies remain safe when the condition changes.\\
91 & Forty edit sequences match the full-calculation oracle.\\
92 & Sorting is numeric and stable for equal keys.\\
93 & Sorting remains a single undoable edit.\\
94 & Descending sort preserves equal-key row order.\\
95 & Search wraps and is case insensitive.\\
96 & Filtering matches display text without changing the workbook.\\
97 & Currency display does not change the stored value.\\
98 & Percent display multiplies only the displayed value.\\
99 & Fixed formatting rounds negative ties away from zero.\\
100 & Cell formats and column widths survive native save and reopen.\\
101 & Formatting is undoable without altering cell contents.\\
102 & Printable HTML escapes workbook text and includes print layout.\\
103 & Expanded sheets calculate references in row 1000.\\
104 & Expanded sparse native files round-trip without empty-cell bloat.\\
105 & Incremental invalidation reaches expanded-sheet references.\\
106 & CSV export includes populated rows beyond row 100.\\
107 & Sort carries formatting with the moved row.\\
108 & Row insertion shifts cell formatting with contents.\\
109 & Copying a rectangle carries cell formatting in one undo step.\\
110 & Font choices survive native serialization.\\
111 & Fill copies formats with values in one undo step.\\
112 & Existing Phase 2 example opens without new style fields.\\
113 & ODS values and formula results round-trip across sheets.\\
114 & ODS text preserves Unicode XML punctuation and literal leading zeroes.\\
115 & ODS basic number formats and bold style round-trip.\\
116 & ODS column widths and alignment round-trip.\\
117 & OpenFormula conversion preserves string delimiters.\\
118 & OpenFormula conversion preserves anchored quoted references.\\
119 & ODS XML parsing resolves namespace URIs rather than literal prefixes.\\
120 & ODS rejects DTD and external entities.\\
121 & ODS rejects malformed XML.\\
122 & ODS rejects merged cells without modifying a workbook.\\
123 & ODS rejects unsupported formulas rather than importing wrong results.\\
124 & ODS package has the spreadsheet MIME type.\\
125 & Packaged ODS opens through the actual ZIP and XML reader.\\
126 & Apache OpenOffice corpus rejects unsupported content explicitly: sampledata.ods.\\
127 & Apache OpenOffice corpus rejects unsupported content explicitly: SubtotalsSampleFile.ods.\\
128 & Format command updates display without changing the formula source.\\
129 & Font command persists a cell font choice.\\
130 & Sort command moves selected rows and their formats.\\
131 & Undo command restores the original row order.\\
132 & Find command selects matching cell text.\\
133 & Filter command records a temporary view without changing contents.\\
134 & Expanded-grid commands edit and calculate cell Z1000.\\
135 & Expanded-grid navigation stops at the actual sheet boundary.\\
136 & =COUNTIF(A1:A3,">15") returns 2.\\
137 & =SUMIF(B1:B3,"apple",A1:A3) returns 40.\\
138 & =AVERAGEIF(A1:A3,">10") returns 25.\\
139 & =SUMPRODUCT(A1:A3,A1:A3) returns 1400.\\
140 & =SUMPRODUCT(A1:A3,A1:A2) returns \#VALUE!.\\
141 & =INDEX(A1:B3,2,2) returns pear.\\
142 & =INDEX(A1:B3,0,1) returns \#REF!.\\
143 & =MATCH("PEAR",B1:B3,0) returns 2.\\
144 & =MATCH(999,A1:A3,0) returns \#N/A.\\
145 & =MATCH(10,A1:A3,1) returns \#N/A.\\
146 & =VLOOKUP(20,A1:B3,2,FALSE) returns pear.\\
147 & =VLOOKUP(20,A1:B3,3,FALSE) returns \#REF!.\\
148 & =LEN("$\lambda$x") returns 2.\\
149 & =LEFT("abcd",2) returns ab.\\
150 & =RIGHT("abcd",2) returns cd.\\
151 & =MID("abcd",2,2) returns bc.\\
152 & =LEFT("abcd",-1) returns \#VALUE!.\\
153 & =UPPER("apple") returns APPLE.\\
154 & =LOWER("PEAR") returns pear.\\
155 & =TRIM("  one   two  ") returns one two.\\
156 & =CONCATENATE("Total ",A1) returns Total 10.\\
157 & =SUBSTITUTE("one two one","one","1") returns 1 two 1.\\
158 & =DATE(1900,1,1) returns 2.\\
159 & =DATE(1900,2,29) returns \#NUM!.\\
160 & =YEAR(DATE(2024,2,29)) returns 2024.\\
161 & =MONTH(DATE(2024,2,29)) returns 2.\\
162 & =DAY(DATE(2024,2,29)) returns 29.\\
163 & =YEAR(DATE(9999,12,31)) returns 9999.\\
164 & =DAY(60) returns 28.\\
165 & =YEAR(1) returns \#NUM!.\\
166 & =SUMIF(A1:A3,">99") returns 0.\\
167 & =AVERAGEIF(A1:A3,">99") returns \#DIV/0!.\\
168 & =NO\_SUCH\_FUNCTION(1) returns \#NAME?.\\
169 & =SUM(revenue) returns 60.\\
170 & =Revenue+1 returns \#VALUE!.\\
171 & =Rate*2 returns 20.\\
172 & Named ranges reject case-insensitive duplicates.\\
173 & Named ranges reject built-in names.\\
174 & Named-range dependencies invalidate cached sums.\\
175 & =SUM(Revenue) returns 60.\\
176 & Named ranges track inserted rows.\\
177 & =SUM(Revenue) returns 60.\\
178 & =DOUBLE(SUM(Revenue)) returns 120.\\
179 & =DOUBLE(1,2) returns \#VALUE!.\\
180 & =RECUR(1) returns \#LIMIT!.\\
181 & =DOUBLE(9) returns 18.\\
182 & Function packs reject unrelated objects.\\
183 & Validation accepts in-range values.\\
184 & Validation rejects out-of-range values.\\
185 & Validation rejects bad pasted data atomically.\\
186 & Validation permits blank cells.\\
187 & Whole-number validation rejects fractions.\\
188 & List validation accepts choices and rejects other text.\\
189 & Conditional colors reflect evaluated numeric values.\\
190 & Conditional formats appear in printable output.\\
191 & Charts read a live source series.\\
192 & Chart export escapes labels and creates SVG.\\
193 & Charts update after source edits.\\
194 & Images accept a complete PNG.\\
195 & Images reject truncated PNG data.\\
196 & Images reject corrupted PNG chunks.\\
197 & Images reject attribute-breaking characters.\\
198 & Printable HTML declares a document language.\\
199 & Native version 4 preserves all workbook objects exactly.\\
200 & Removing an object is undoable.\\
201 & ODS refuses to silently drop native workbook objects.\\
202 & Pivot tables group repeated labels and sum values.\\
203 & Pivot creation is one undo transaction.\\
204 & Pivot refresh reads changed source values.\\
205 & Pivot averages retain exact numeric results.\\
206 & Pivot definitions survive native reopen and refresh.\\
207 & A sparse working sheet stores only populated cells.\\
208 & Clearing a sparse cell releases its entry.\\
209 & The display cache stores no blank-cell results.\\
210 & =SUM(Revenue) returns 60.\\
211 & Cut moves complete named ranges with their cells.\\
212 & Cut rejects partial named-range moves.\\
213 & Pivot sources follow sheet renames.\\
214 & Pivot sources follow inserted source columns.\\
215 & Deleting a pivot key column is rejected explicitly.\\
216 & Custom functions recalculate after named inputs change.\\
217 & CSV replacement obeys validation rules.\\
218 & Date formatting preserves serial values and displays ISO dates.\\
219 & Date formatting survives native reopen.\\
220 & Thirty-two large sparse sheets retain only thirty-two populated entries.\\
221 & Line charts render a marker for a single data point.\\
222 & Charts over a range with no numbers say so.\\
223 & Named-range menu saves the selected rectangle.\\
224 & Custom-function menu defines a callable function.\\
225 & Rejected cell validation preserves the workbook and draft.\\
226 & Conditional-format menu paints matching cells.\\
227 & Chart menu stores a source range and renders through canvas commands.\\
228 & Undo removes a newly inserted chart.\\
229 & Object removal uses one reversible workbook edit.\\
230 & Function-pack menus round-trip declarative definitions.\\
231 & PNG preflight accepts a small valid header.\\
232 & PNG preflight rejects non-image input.\\
233 & PNG preflight rejects excessive decoded dimensions.\\
234 & Pivot menu creates and activates a summary sheet.\\
235 & An unchanged pivot refresh does not add undo history.\\
236 & Finder open requests wait until the interface is ready.\\
237 & Queued Finder requests open a native workbook after startup.\\
238 & Dropped paths preserve Unicode spaces and Tcl metacharacters.\\
239 & Cancel keeps pending edits when another file is dropped.\\
240 & Cancelled drops do not overwrite the original file.\\
241 & Discard opens a file dropped onto a running app.\\
242 & Invalid dropped workbooks leave the current workbook intact.\\
243 & Missing dropped files leave the current workbook intact.\\
244 & Multiple-file drops are rejected without silently replacing workbooks.\\
245 & Unsupported file types receive a clear rejection.\\
246 & Dropped CSV files import as new unsaved workbooks.\\
247 & Dropped CSV formula-looking fields remain literal text.\\
248 & Cancelling Save As cancels the document switch.\\
249 & Dropped ODS files import without becoming the native Save target.\\
250 & Document events cannot reenter an in-progress open dialog.\\
251 & Deferred document events resume after the current operation.\\
252 & Save-before-open persists edits before installing the dropped workbook.\\
% END ADVANCED CATALOGUE
\bottomrule\end{longtable}

\clearpage
\section{Formal verification}
\clearpage
\section{Staged road map}
Each phase depends on the preceding phase and should be completed and validated
before the next begins. Phase 1 establishes a deliberately narrow spreadsheet;
OpenOffice feature parity is a long-term objective. Estimates and completion dates
will depend on compatibility choices and measured implementation work.

\subsection*{Phase 1: minimal single-sheet spreadsheet}
\textbf{Scope.} Text and exact numbers; direct in-cell editing; A1 references; arithmetic and
SUM; cycle/error detection; undo/redo; single-cell clipboard operations; native
open/save; basic menus and shortcuts; Lean-to-C native build; this manual and the
61-test baseline. No images, charts, macros, formatting system, or office-file
compatibility.

\textbf{Current outcome.} Implemented and compiled for Apple-silicon macOS; 61
native tests pass. A separate GUI smoke test and launch diagnostics are supplied.
The user confirmed responsive navigation after the 0.1.4 rendering fix.
That implementation is retained in Phase 2.

\textbf{Exit gate.} Validate normal interactive startup, cell editing, navigation,
recalculation, clipboard, undo/redo, save/reopen, conflict reporting, and cancellation
of unsaved-change prompts in a logged-in desktop session. Retain a reproducible
native build and an up-to-date PDF. Do not mark this gate complete on native unit
tests alone.

\subsection*{Phase 2: everyday workbook operations}
\textbf{Scope.} Multiple named sheets; cross-sheet references; relative, absolute,
and mixed references; rectangular selection/copy/paste; fill; row/column insertion
and deletion; CSV import/export with explicit delimiter and encoding rules; more
functions such as MIN, MAX, AVERAGE, COUNT and IF.

\textbf{Current outcome.} Implemented in 0.5.1 with the rules specified before
implementation. The new format, migration, and engine/handler checks pass.
Phase 2 desktop acceptance remains pending because this sandbox cannot start Tk.

\textbf{Exit gate.} Specify reference adjustments before implementation; regression
tests must cover structural edits, quoted CSV fields, cross-sheet cycles and undo
across sheets. Version the native format and test migration from Phase 1 files.

\subsection*{Phase 3: usable scale and presentation}
\textbf{Scope.} Dependency tracking and incremental recalculation; larger sparse
sheets; responsive rendering; search, sort and filtering; column widths, fonts,
alignment, numeric display formats and print layout. Separate stored values from
display formatting and agree on decimal/date/rounding semantics before expanding
the numeric model.

\textbf{Current outcome.} Implemented to the bounded 1,000-row scope in 0.5.1,
with sparse working storage and sparse display caches, incremental invalidation,
search/sort/filter, fonts/formats, explicit Gregorian date semantics, and printable
HTML. The named performance fixture meets its time and memory budgets; desktop
painting and print-dialog acceptance remain pending.

\textbf{Exit gate.} Publish measured memory and recalculation budgets using named
fixtures; test edits against a full-recalculation oracle; verify sort stability,
undo, error propagation and print output. Keep calculations independent of the UI.

\clearpage
\subsection*{Phase 4: OpenDocument interoperability}
\textbf{Scope.} Read and write a declared subset of ODS/OpenFormula, starting with
values, formulas and basic styles. Add explicit diagnostics for unsupported
features and a policy for preserving them. Treat XLSX as a separate later adapter,
not an automatic consequence of ODS support.

\textbf{Current outcome.} Initial ODS subset implemented and generated round trips
tested in 0.5.1. A pinned Apache OpenOffice corpus verifies explicit rejection
of unsupported validation and styling. The feature matrix and source hashes are
recorded under \code{lean_calc/fixtures/ods}. Live OpenOffice comparison is blocked
by this session: LaunchServices is unavailable and mounting the official installer
returns ``Device not configured. This exit gate is not complete.

\textbf{Exit gate.} Use a versioned corpus of real and generated workbooks; compare
results and round trips with Apache OpenOffice Calc. Publish the supported feature
matrix and distinguish exact round trips from intentional conversions. Never
silently drop unsupported content when saving an imported workbook.

\subsection*{Phase 5: advanced Calc capabilities}
\textbf{Scope.} Add advanced functions, named ranges, validation, conditional
formatting, charts, images, pivot tables, and extensibility in separately testable
increments. Prioritize actual workbook needs. Consider macros only after defining
a clear execution and permission model.

\textbf{Current outcome.} Advanced functions, names, validation, conditional
formatting, charts, PNG images, pivots and declarative function packs have models,
native persistence, menus, fixtures and regression tests. Desktop appearance and
interaction acceptance are still blocked. Executable macros are deferred until
an execution/permission model is explicitly designed.

\textbf{Exit gate.} Each feature needs a data model, compatibility fixtures,
regression coverage, UI acceptance and updated documentation. Maintain a clear
boundary between the calculation engine and optional presentation features.

\subsection*{Longer-term continuation: other OpenOffice applications}
After the spreadsheet has reached an agreed usable scope, plan separate phases
for a Writer-like document editor and then presentation/drawing tools. Reuse the
Lean build, file handling, UI infrastructure and testing approach where practical.
The current release implements bounded Phase 3 and Phase 5 features plus an ODS subset; it does not port
OpenOffice's existing codebase or replace the entire suite.

\subsection*{Documentation and acceptance at every phase}
Compile \code{LeanCalc.tex} to \code{LeanCalc.pdf} on every LeanCalc work turn.
Keep menus and shortcuts synchronized with implementation. Preserve the original
61 test identifiers and append new regression tests as features grow. Keep the
formal-verification section empty until a later explicit instruction opens that
work; no proof search or ranking is undertaken in this version.
\end{document}
